\chapter{3}

\section[«Roccia del Culto», Mailand (IT)]{\emph{«Roccia del Culto»}, Mailand (IT)}

Den Blick schliesslich an sie geheftet, sah er ihr Gesicht. Doch selbst wenn sie ihm zuwinkte und ihr Mund zu einem Lächeln geformt war, 
die Entfernung konnte ihn nicht täuschen. Es wirkte gequält. 

Louis suchte den Platz hektisch nach Valentina ab. Sie schaute hoch und klopfte mit der flachen Hand auf einen leeren Platz neben sich. Er ging hin, legte den Rucksack vor sich nieder, setzte sich und stellte den Gitarrenkoffer ab. Die Musik dominierte die Szene. Unerwartet tauchten vor Louis eine Bierflasche und Francos freundliches Gesicht auf.

«Für dich.»

Er schrie die Worte in den musikalischen Schwung hinein, prostete Louis zu, blieb ein paar Takte neben ihm stehen und versuchte, ihn zum
Mittanzen zu animieren. Louis blieb sitzen. Und Franco schien zu verstehen. Das kühle Bier tat gut.

Louis starrte in die Flammen. Bis sich sein Blick in Richtung Giorgia verlor und ihm Valentina unerwartet über den Rücken strich. Er sah sie stumm an, kramte verzweifelt nach Leichtigkeit, nach Verbindung, nach irgendwelchen Möglichkeiten, ins Gespräch zu kommen. Nichts wollte ihm einfallen. Shit. Valentina nahm ihre Hand zurück.

«Ist ok, brauchst nichts zu sagen.»

Jemand legte eifrig Holz nach und kurz darauf erhellte sich die Umgebung. Es würde nicht mehr lange dauern, bis eine Streife der «Polizia Nationale» eintraf, wie so oft, \emph{«Trockenheit -- Feuer entfachen verboten»} würde der Ordnungshüter sagen. Sie würden die verlangte Busse zusammenlegen und mit ernsten Minen versichern, das Feuer augenblicklich zu löschen. Heute kam keiner. Und zum ersten Mal hätte sich Louis über den Auftritt der Polizisten gefreut.

Verstohlen schaute er um sich. Wie befremdend er ihre Ausgelassenheit wahrnahm. Sehnsüchtige, alkoholtrunkene Blicke huschten von da nach dort und manchmal zurück. Unterdessen war es wieder ruhiger geworden, wie Musik und Gesang vorzugeben schienen. Bis \emph{«Nera»} erklang, in ihrer Runde omnipräsent, den ganzen Sommer schon. Alle wussten, es war Bastianos Wahl. Der Song war längst zum untrüglichen Signal geworden. Er war wieder auf Eroberungstour.

\qrblock{Irama \emph{«Nera»}}{media/QR-Codes/057_Nera_Irama_Spotify.png}

Da stand Bastiano unterdessen, schwankend, zusammen mit einer Frau, die Louis noch nie gesehen hatte. An einen Baum gelehnt schien er sich
unter Küssen aufzulösen. Bald zogen sich die beiden hinter die Büsche zurück.

Louis wandte seinen Blick entschieden in eine andere Richtung. Sein Schmerz dehnte sich plötzlich, schnell und beissend in ihm aus, als würden ihm Rhythmus und Text des Songs geradezu intravenös verabreicht. Wie hemmungslos sich Bastiano einschmeichelte. Wie direkt. Wie banal. Selbst Giorgia hatte ihm beigepflichtet, als sie sich «Nera» zum ersten Mal angehört hatten. Von billiger Anmache, hatte sie gesprochen. Einem Frauenbild, das sie anwidere. Doch auch sie liess ihn gewähren. Wie die andern. Keiner sei perfekt, sagte sie entschuldigend und nannte ihn weiterhin \emph{ihren Bruder}. Louis verstand sie nicht.

Es war diese bekannte Stimme der Missgunst, die immer wieder aufflammte. Heftige Ablehnung und eine Spur von Bewunderung kämpften gegeneinander. Bastiano war ein elender Bastard. Er hatte etwas, das Louis fehlte. Der Typ schien genau zu wissen, wann er welche Knöpfe bei den Frauen drücken musste und selbst, wo sie sie versteckt hielten und bekam, was er wollte. Schamlos. Er spielte den König der Verführung. Frauen fielen reihenweise auf ihn rein.
Mit jedem Gedanken steigerte sich Louis` Wut. Und wieder kam diese Angst hoch, Giorgia könnte zu Bastiano überspringen. Auch sie würde er einnehmen, ausnutzen und später fallenlassen. Selbst ihre Trennung hatte nichts an seiner Entrüstung geändert.

\sectionbreak

Bastiano war inzwischen aus dem Dunkel aufgetaucht. Er warf mehrere \emph{Ciaos} in die Runde. Dann ging er. Die junge, unbekannte Frau war bereits kurz zuvor verschwunden.
Louis beobachtete Giorgia nun vorwiegend aus den Augenwinkeln. Wann würde sie zu ihm schauen, wenigstens ein zweites Mal in dieser Nacht? Wohin sie nur dauernd starrte? Er rätselte, verengte die Augen, konnte aber ihren Seitenblick nicht deuten. Da gab es nur die Schatten einzelner Bäume neben und hintereinander. Sonst nichts. Es war sonderbar.

\sectionbreak

Hatte es wirklich nur an seiner verschlossenen Art gelegen?
Oder war es \emph{der} Vorwand, weil sie bereits einen Andern hatte?
Die Frage kam das erste Mal in dieser Klarheit und mit einer Kälte, die ihn
frösteln liess. Er verkrampfte sich. Daran hatte er nicht gedacht. Wie
dumm er war. Ein anderer. Das würde alles erklären. Die Zurückweisung.
Auch die Funkstille nach Louis’ Bitte um eine letzte Aussprache. Und er
entschied, jetzt, in dieser Nacht, wollte er die Antwort. Sein Atem ging
schneller.
Es sollte seine letzte Frage an Giorgia sein. 
Er würde es wissen. Es würde endgültig sein. Eine abgeschlossene Sache, 
wenn auch, so oder so, ein einziges Elend. Es machte keinen Unterschied
mehr.

Als wären Louis’ Gedanken auf Giorgia übergesprungen, stand sie umständlich auf. Ob er jetzt auf sie zugehen sollte? 
Was hatte sie vor? Giorgia zielte an Francos Beinen vorbei. Dann sagte sie etwas zu Maria. Die schien ihr zu antworten. 
Giorgia schüttelte den Kopf. Maria schaute ihr nach. Sie sah besorgt aus. Giorgia verschwand schwankend zwischen den Bäumen. Dabei lachte sie heiser.
So hatte er sie kaum je gesehen. Sie, die zwar mittrank, sich aber auch klare Grenzen setzte. Sie halte ihre Sinne gerne beisammen. Er hatte sie dafür bewundert. Wie überlegt sie war. Und nun dieser Auftritt. Es war offensichtlich, sie machte einen vollkommen zugedröhnten Eindruck.
Wollte sie ihn provozieren? Sie musste gewusst haben, dass er heute da sein würde.

Louis fühlte sich von Valentina beobachtet. Sie schien seine Anspannung wahrzunehmen und ihm zuzusehen, wie er Giorgias Bewegungen verfolgte. Er blieb noch eine Weile unbeweglich sitzen.
Francesco Gabbanis \emph{«Occidentali’s Karma»} klang unterdessen überlaut zwischen den Beinen der Tanzenden hindurch und hielt die Nacht
am Leben.

\qrblock{Francesco Gabbani \emph{«Occidentali’s Karma»}}{media/QR-Codes/009_Occidentali's_Karma_(Radio_Edit)_Francesco_Gabbani_Spotify.png}

Dann der Wechsel. Abrupt. Die Emotionen wandelten sich. 
\emph{«The Feeling Begins»} löste den eingängigen Elektrosound ab.

\qrblock{Peter Gabriel \emph{«The Feeling Begins»}}{media/QR-Codes/003_The_Feeling_Begins_Peter_Gabriel_Spotify.png}

Die Leichtigkeit kippte. Louis war, als lege sich eine schwere Decke über den Ort. Einige wiegten sich hin und her, als wären sie in Trance und würden es ewig bleiben. Andere schienen die Wirkung eingeworfener Pillen still auszukosten. Die meisten rückten näher ans Feuer. Heute war einiges im Umlauf.

Jetzt.
Louis holte Luft.
